\subsection{珍稀物种响应的归因能力}
问题二把公报初值、通道阻隔、放流补偿和上游食源纳入同一响应框架后，连通性约束与工程补偿的异质作用得以区分。高阻隔会使江豚和长江鲟末值降至 0.901 和 0.856，无放流时长江鲟降至 0.519；配合 $\pm 20\%$ 参数敏感性结果看，江豚对自身摄食与转化参数更敏感，长江鲟则对放流强度更敏感。

这一部分的限制主要来自状态结构的简化。当前模型尚未展开分龄、繁殖季节和分阶段放流调度，也把江豚、长江鲟与入侵者对中层猎物的半饱和响应压缩在同一尺度上。按这一级别的简化，它更适合比较不同机制的相对强弱，而不直接外推为精细的长期种群预测。

